% !PRINTED 19
\weberpage{19}
Dies folgt aber unmittelbar aus

\weq{7}{Q_1 a^{-1} Q_1 b Q_1 a = Q_1 a^{-1} b a,}

weil zugleich mit $b$ auch $a^{-1} b a$ in $Q$ vorkommt (weil $Q$ ein
Normaltheiler von $P$ ist). Und dass auch die Bestimmung der
Indices richtig ist, zeigt die Abzählung. Denn die Grade von
$P/Q_1$ und $Q/Q_1$ sind $\nu \nu_1$ und $\nu_1$, also der Index $\nu$.

Daneben besteht folgender Satz:

\gesperrt{2. Ist $Q_1$ ein Normaltheiler von $P$ vom Grade $q_1$,
und hat die Gruppe $P/Q_1$ selbst einen Theiler $M$
vom Grade $m$, so hat $P$ einen Theiler $Q$ vom
Grade $q = m q_1$. Zugleich ist $Q_1$ ein Normaltheiler
 von $Q$, und es ist $M = Q/Q_1$. Ist $M$ Normaltheiler
 von $P/Q_1$, so ist auch $Q$ Normaltheiler
von $P$.}

Es sei $\mu$ der Index von $Q_1$ in Bezug auf $P$, und $P$ sei in
die Nebengruppen zerlegt:

\weq{8}{P = Q_1 + Q_1 e_1 + Q_1 e_2 \cdots + Q_1 e_{\mu - 1},}

so dass $e_1, e_2 \ldots e_{\mu - 1}$ passend gewählte Elemente in $P$ und

\weq{9}{Q_1,\, Q_1 e_1,\, Q_1 e_2 \ldots Q_1 e_{\mu - 1}}

die Elemente von $P/Q_1$ sind. Wenn nun in der Gruppe $P/Q_1$
ein Theiler $M$ enthalten ist, so muss dieser aus einem Theile
der Elemente (9) bestehen, etwa aus

\weq{10}{Q_1,\, Q_1 b_1,\, Q_1 b_2 \ldots Q_1 b_{\nu_1 - 1},}

und wenn dies System eine Gruppe bilden soll, so muss für
irgend zwei Elemente $b_i, b_k$ das Product

\weq{11}{Q_1 b_i Q_1 b_k = Q_1 b_i b_k}

wieder in (10) enthalten, also etwa $= Q_1 b_h$ sein. Wenn also
das System $1, b_1, b_2 \ldots b_{\nu_1 - 1}$ mit $B$ bezeichnet wird, so lässt
sich nach der Methode der Composition der Theile die Relation
 (11) so schreiben:

\weq{12}{Q_1 B Q_1 B = Q_1 B,}

und so besagt sie nach \S.\,4, 1., dass die in

\weq{13}{Q = Q_1 B}

enthaltenen Elemente von $P$ eine Gruppe bilden, die die Gruppe
$Q_1$ als Theiler enthält und selbst ein Theiler von $P$ ist. Dass
$Q_1$ Normaltheiler von $Q$ ist, folgt aus \S.\,4, 5., weil $Q_1$ als Normal-
%CONT>>
